\documentclass[10pt,a4paper]{amsart}
\usepackage[margin=19mm]{geometry}
\usepackage{amsmath,amssymb,mathtools}
\usepackage[scr=rsfs,cal=euler]{mathalfa}
\usepackage{xfrac}
\usepackage[unicode,hidelinks,bookmarks=false]{hyperref}
\newcommand{\ie}{i.e.\ }
\newcommand{\cf}{cf.\ }
\newcommand{\resp}{respectively\ }
\newcommand{\Subsection}{\subsection}
\providecommand{\nobreakdash}{\nobreak\hbox{-}}
\setlength{\emergencystretch}{2em}
\multlinegap0pt
\usepackage[shortlabels]{enumitem}
\usepackage{amssymb}
\usepackage{mathtools}
\usepackage{float}
\usepackage{algorithm}\usepackage{algpseudocode}
\usepackage{xcolor}
\newtheorem{lemma}{Lemma}
\newtheorem{theorem}{Theorem}
\title{Accelerated Convergence of a Second-Order Dynamical System and its Application to Splitting Algorithms for Comonotone Inclusions}
\author{Yan Tang, Jun Dong}
\date{Source: arXiv:2609.02479v1 (CC BY 4.0)}
\begin{document}
\maketitle

\noindent\emph{Source and licence.} Accelerated Convergence of a Second-Order Dynamical System and its Application to Splitting Algorithms for Comonotone Inclusions. Version arXiv:2609.02479v1 (CC BY 4.0), distributed by arXiv under the Creative Commons Attribution 4.0 International licence, http://creativecommons.org/licenses/by/4.0/, as recorded in the arXiv licence metadata for this exact version and shown on the abstract page https://arxiv.org/abs/2609.02479v1. Full licence notice: \texttt{licenses/arxiv-2609.02479-CC-BY-4.0.txt}.\par\smallskip

\noindent\textit{Editorial scope.} Counted unit: the article's theorem thm:pointwise\_magnitudes\_layered (source0.tex lines 727--753). The article's complete original proof of it is delivered with it (source0.tex lines 755--845, one \texttt{proof} environment of 91 lines). No mathematics is written, repaired or re-derived: the statement and the proof are the article's own lines, in the article's own notation, and no retained line is substituted or re-flowed.  Retained uncounted as the source-local context the proof needs: lines 318--324; lines 372--375; lines 554--559; lines 565--585; lines 677--694.  Nothing was omitted from any retained span, so the package carries no omission record.  No retained line contains a citation, so the wrapper carries no bibliography and no bibliography entry.  No retained line contains a figure environment, an image-inclusion command or drawing code, so no image is delivered and the package declares no asset dependency.  Editorial formatting only, in addition: an AMS \texttt{amsart} preamble that restates the article's own declarations and macros used by the retained lines, namely the environments \texttt{lemma}, \texttt{theorem} on separate counters, together with the standard packages those lines need.  The retained environments print in this document as Lemma 1, Theorem 1 in order of appearance, whereas the article prints the same items with the numbering of its own full document.  The complete native source of the article as distributed in the arXiv e-print for this exact version is bundled unchanged as \texttt{sources/209204/source0.tex}.
\begin{lemma}
\label{lem:weighted_decay_criterion}
Let $r:[t_0,+\infty)\to(0,+\infty)$ be locally integrable and satisfy
$\int_{t_0}^{+\infty}r(t)dt=+\infty$. Let $F:[t_0,+\infty)\to\mathcal{H}$ be locally absolutely continuous. Assume that $\int_{t_0}^{+\infty} r(t)\|F(t)\|_{\mathcal{H}}^2dt<+\infty$ and that, for some $t_1\ge t_0$ and $L>0$, $\|\dot F(t)\|_{\mathcal{H}}\le Lr(t)$ 
for almost every $t\ge t_1$. Then
$\lim_{t\to+\infty}\|F(t)\|_{\mathcal{H}}=0$.
\end{lemma}

\begin{equation}
\label{eq:SHDS}
\ddot{x}(t) + \alpha(t)\dot{x}(t) + \beta(t)\frac{d}{dt}T_{\lambda(t), \gamma(t)}(x(t)) + b(t)T_{\lambda(t), \gamma(t)}(x(t)) = 0,
\end{equation}

\begin{lemma}
\label{prop:omega_uniform_positivity}
Under Assumption {\bf (H1)}, the function 
$\Omega(t) := \Omega(\gamma(t)) = \frac{4\nu(\rho + \gamma(t)) - \gamma^2(t)}{4\gamma(t)(\nu + \rho)}$ 
is bounded away from zero on $[t_0, +\infty)$. That is, there exists a constant $\underline{\Omega} > 0$ such that $\Omega(t) \ge \underline{\Omega}$ for all $t \ge t_0$.
\end{lemma}

\begin{lemma}
\label{lem:sharp_sensitivity_final}
Under Assumption {\rm {\bf (H1)}}, let $A: \mathcal{H} \to 2^{\mathcal{H}}$ be a maximal $\rho$-comonotone operator and $B: \mathcal{H} \to \mathcal{H}$ be a $\nu$-cocoercive operator with $\nu > 0$. For any $\lambda_1, \lambda_2 > 0$ and $\gamma_1, \gamma_2 \in \left( \max(0, -2\rho), \, 2\nu \right)$, let $\Lambda(t) := \lambda(t) \Omega(t) > 0$ be the cocoercivity constant of $T_{\lambda(t), \gamma(t)}$ with $\Omega(t) := \frac{4\nu \left( \rho + \gamma(t) \right) - \gamma^2(t)}{4\gamma(t) \left( \nu + \rho \right)}$. For any $\bar{x} \in \operatorname{zer}(A+B)$ and $x, y, z \in \mathcal{H}$, the following assertions hold:

{\rm (i)} The resolvent operator satisfies the following inequality:
\begin{equation}\label{eq:resolvent_sensitivity_standard}
\| J_{\gamma_1 A} z - J_{\gamma_2 A} z \|_{\mathcal{H}} \leq \frac{|\gamma_1 - \gamma_2|}{\gamma_1} \| z - J_{\gamma_1 A} z \|_{\mathcal{H}}.
\end{equation}

{\rm (ii)} The difference of the operators satisfies:
\begin{align}
\|\lambda_{1} T_{\lambda_{1}, \gamma_{1}}(x) - \lambda_{2} T_{\lambda_{2}, \gamma_{2}}(y)\|_{\mathcal{H}} \leq \;& 4 \left( \frac{\lambda_1}{\Lambda_1} + 4 \right)\frac{\nu}{\gamma_1} \|x - y\|_{\mathcal{H}} + \frac{2 |\gamma_{1} - \gamma_{2}|}{\gamma_{1}} \| x - \bar{x} \|_{\mathcal{H}} \nonumber \\
& + \frac{8\nu |\gamma_{1} - \gamma_{2}|}{\gamma_{1}} \|B(x)\|_{\mathcal{H}}, \label{eq:asymptotic_consistent_form}
\end{align}
where $\Lambda_1$ denotes $\Lambda(t_1)$ evaluated at parameters $(\lambda_1, \gamma_1)$.

{\rm (iii)} For any strong global solution $x(t)$ of system \eqref{eq:SHDS}, the derivative of the operator satisfies for almost every $t \ge t_0$:
\begin{equation}\label{eq:final_CL_derivative}
\left\| \frac{d}{dt} \left( \lambda(t) T_{\lambda(t), \gamma(t)}(x(t)) \right) \right\|_{\mathcal{H}} \leq 4 \left( \frac{1}{\underline{\Omega}} + 4 \right) \frac{\nu}{\gamma(t)} \|\dot{x}(t)\|_{\mathcal{H}} + \frac{2 |\dot{\gamma}(t)|}{\gamma(t)} \| x(t) - \bar{x} \|_{\mathcal{H}} + \frac{8\nu |\dot{\gamma}(t)|}{\gamma(t)} \| B(x(t)) \|_{\mathcal{H}}.
\end{equation}
\end{lemma}

\begin{theorem}
\label{thm:m_weight_bounds}
Assume that Assumptions {\rm {\bf (H1)-(H3)}} hold and let $x(\cdot):[t_0, +\infty) \to \mathcal{H}$ be a solution trajectory of system \eqref{eq:SHDS}. Suppose that the following conditions are satisfied for almost every $t \geq t_0$:
\begin{equation*}
\begin{aligned}
    &\text{(C1)} \quad m(t) \dot{a}(t) + a(t) \left( a(t) + \dot{m}(t) - m(t) \alpha(t) \right) + c(t) = 0; \quad \text{(C2)} \quad \mathcal{W}_v(t) > 0, \, \mathcal{W}_T(t) > 0;\\ \quad 
    &\text{(C3)} \quad a(t)\dot{a}(t) + \frac{1}{2}\dot{c}(t) \le 0; \quad
    \text{(C4)} \quad \inf_{t \ge t_0} c(t) = \underline{c} > 0.
\end{aligned}
\end{equation*}
Then, the functional $\mathcal{E}(t)$ is non-increasing on $[t_0, +\infty)$ and the following properties hold:
\begin{enumerate}[label=(\roman*)]
    \item $\sup_{t \geq t_0} \|x(t) - \bar{x}\|_{\mathcal{H}}< +\infty$;
    \item $\sup_{t \geq t_0} \left\| a(t) \left( x(t) - \bar{x} \right) + m(t) \left( \dot{x}(t) + \beta(t) T_{\lambda(t), \gamma(t)}(x(t)) \right) \right\|_{\mathcal{H}} < +\infty$;
    \item $\int_{t_0}^{+\infty} \mathcal{W}_T(t) \| T_{\lambda(t), \gamma(t)}(x(t)) \|_{\mathcal{H}}^2 dt < +\infty$;
    \item $\int_{t_0}^{+\infty} \mathcal{W}_v(t) \| \dot{x}(t) + \beta(t) T_{\lambda(t), \gamma(t)}(x(t)) \|_{\mathcal{H}}^2 dt < +\infty$.
\end{enumerate}
\end{theorem}

\begin{theorem}
\label{thm:pointwise_magnitudes_layered}
Suppose Assumptions {\bf (H1)--(H3)} and conditions {\rm (C1)--(C4)} hold. Let $x(t)$ be a solution trajectory of system \eqref{eq:SHDS}. As $t \to +\infty$, the following assertions hold:

\smallskip
\noindent \textbf{(1) Asymptotic Magnitudes.} If {\rm (C5)} $\bar{a} := \sup_{t \ge t_0} a(t) < +\infty$ and $\sup_{t \geq t_0} \frac{m(t) \beta(t)}{\lambda(t)} < +\infty$ are satisfied, then:
\begin{enumerate}[label=(\roman*), nosep]
    \item $\|T_{\lambda(t), \gamma(t)}(x(t))\|_{\mathcal{H}} = \mathcal{O}( \frac{1}{\lambda(t)})$;
    \item $\|\dot{x}(t)\|_{\mathcal{H}} = \mathcal{O}( \frac{1}{m(t)} )$.
\end{enumerate}

\smallskip
\noindent \textbf{(2) Operator Derivative Magnitude.} If, in addition to {\rm (C5)}, parameters fulfill {\rm (C6)} $0<\underline{m}' \leq \dot{m}(t) \leq \bar{m}'$, $m(t) \geq \dot{m}(t)$, $\underline{\gamma} := \inf_{t \geq t_0} \gamma(t) > 0$; and {\rm (C7)} $\frac{|\dot{\gamma}(t)|}{\gamma(t)}, \frac{|\dot{\lambda}(t)|}{\lambda(t)} = \mathcal{O}( \frac{\dot{m}(t)}{m(t)} )$, then:
\begin{enumerate}[label=(\roman*), resume, nosep]
    \item $\| \frac{d}{dt} T_{\lambda(t), \gamma(t)}(x(t)) \|_{\mathcal{H}} = \mathcal{O}\left( \frac{\dot{m}(t)}{m(t) \lambda(t)} \right)$.
\end{enumerate}

\smallskip
\noindent \textbf{(3) Accelerated Trajectory and Residual Rates.} Building upon conditions in \textbf{(1)} and \textbf{(2)}, if parameters satisfy {\rm (C8)} $\alpha(t) = \mathcal{O}( \frac{\dot{m}(t)}{m(t)} )$, $\beta(t) = \mathcal{O}( \frac{m(t)}{\dot{m}(t)} )$, $b(t) = \mathcal{O}(1)$, $b^2(t) \leq C_b \frac{\dot{m}(t) \mathcal{W}_T(t)}{m^3(t)}$; {\rm (C9)} $\lambda(t) \ge C_{\lambda} m^2(t)$; {\rm (C10)} $\mathcal{W}_v(t) \ge C_v m(t)\dot{m}(t)$, $\mathcal{W}_T(t) \ge C_T m^2(t)\dot{m}(t)$ for some $C_b, C_{\lambda}, C_v, C_T > 0$; and {\rm (C11)} $\liminf_{t \to +\infty} \frac{m(t) \mathcal{W}_T(t)}{\dot{m}(t) \lambda^2(t)}:=w > 0$, then:
\begin{enumerate}[label=(\roman*), resume, nosep, itemsep=2pt]
    \item $\|\ddot{x}(t)\|_{\mathcal{H}} = \mathcal{O}\left( \frac{\dot{m}(t)}{m^2(t)} \right)$;
    \item $\|\dot{x}(t)\|_{\mathcal{H}} = o\left( \frac{1}{m(t)} \right)$;
    \item $\|T_{\lambda(t), \gamma(t)}(x(t))\|_{\mathcal{H}} = o\left( \frac{1}{\lambda(t)} \right)$;
    \item $\| A_{\gamma(t)} ( x(t) - \gamma(t) Bx(t) ) + Bx(t) \|_{\mathcal{H}} = o\left( \frac{1}{\gamma(t)} \right)$;
    \item $\| \frac{d}{dt} ( A_{\gamma(t)} ( x(t) - \gamma(t) Bx(t) ) + Bx(t) ) \|_{\mathcal{H}} = \mathcal{O}\left( \frac{\dot{m}(t)}{m(t) \gamma(t)} \right) + o\left( \frac{\lambda(t) | \frac{d}{dt} ( \frac{\gamma(t)}{\lambda(t)} ) |}{\gamma^2(t)} \right)$.
\end{enumerate}
\end{theorem}

\begin{proof}
\hfill \\
\noindent \textbf{(1) Basic Asymptotic Magnitudes:} By combining the $\frac{1}{\Lambda(t)}>0$-Lipschitz continuity of $T_{\lambda(t), \gamma(t)}$ with the boundedness of the trajectory $\|x(t) - \bar{x}\|_{\mathcal{H}} \le M_x$ established in Theorem \ref{thm:m_weight_bounds}, assertion (i) follows from the estimate:
\begin{equation}
\label{T__first_rate}
\|T_{\lambda(t), \gamma(t)}(x(t))\|_{\mathcal{H}} \le \frac{1}{\lambda(t) \Omega(t)} \|x(t) - \bar{x}\|_{\mathcal{H}} \le \frac{M_x}{\underline{\Omega} \lambda(t)} = \mathcal{O}\left( \frac{1}{\lambda(t)} \right),
\end{equation}
where $\Omega(t):=\frac{4\nu \left( \rho + \gamma(t) \right) - \gamma^2(t)}{4\gamma(t) \left( \nu + \rho \right)}$ and $\underline{\Omega} := \inf_{t \ge t_0} \Omega(t)> 0$ from Lemma \ref{prop:omega_uniform_positivity}. Utilizing Theorem \ref{thm:m_weight_bounds} (ii), condition (C5), \eqref{T__first_rate} and triangle inequality, we have:
\begin{align*}
\|m(t) \dot{x}(t)\|_{\mathcal{H}} &\le \left\| a(t) \left( x(t) - \bar{x} \right) + m(t) \left( \dot{x}(t) + \beta(t) T_{\lambda(t), \gamma(t)}(x(t)) \right) \right\|_{\mathcal{H}} \nonumber \\
&\quad + a(t)\|x(t) - \bar{x}\|_{\mathcal{H}} + m(t)\beta(t)\|T_{\lambda(t), \gamma(t)}(x(t))\|_{\mathcal{H}} \nonumber \\
&\le \sqrt{2\mathcal{E}(t_0)} + \bar{a} M_x + \left( \frac{M_x}{\underline{\Omega}} \right) \left( \frac{m(t)\beta(t)}{\lambda(t)} \right)\nonumber \\
&< + \infty,
\label{eq:m_isolation_refined}
\end{align*}
which yields $\|\dot{x}(t)\|_{\mathcal{H}} = \mathcal{O}\left( \frac{1}{m(t)} \right)$, thus establishing assertion (ii).

\noindent \textbf{ (2) Operator Derivative Magnitude:}
For assertion (iii), by combining Lemma \ref{lem:sharp_sensitivity_final} (iii), the boundedness of the trajectory, the velocity magnitude from Theorem \ref{thm:pointwise_magnitudes_layered} (ii), and conditions (C6)--(C7), along with the Lipschitz continuity of $B$, it follows that as $t \to +\infty$:

\begin{align}
\left\| \frac{d}{dt}\lambda(t)T_{\lambda(t),\gamma(t)}(x(t)) \right\|_{\mathcal{H}} &\le  \frac{4 \left( \frac{1}{\underline{\Omega}} + 4 \right)\nu}{\gamma(t)} \left\| \dot{x}(t) \right\|_{\mathcal{H}} + \frac{2 |\dot{\gamma}(t)|}{\gamma(t)} \| x(t) - \bar{x} \|_{\mathcal{H}} + \frac{8\nu |\dot{\gamma}(t)|}{\gamma(t)} \| B(x(t)) \|_{\mathcal{H}} \nonumber \\
&\le \frac{4 \left( \frac{1}{\underline{\Omega}} + 4 \right)\nu}{\underline{\gamma} \underline{m}'} \mathcal{O}\left( \frac{\dot{m}(t)}{m(t)} \right) +\frac{2 |\dot{\gamma}(t)|}{\gamma(t)} \| x(t) - \bar{x} \|_{\mathcal{H}} + \frac{8\nu |\dot{\gamma}(t)|}{\gamma(t)} \| B(x(t)) \|_{\mathcal{H}} \nonumber \\
&= \mathcal{O}\left( \frac{\dot{m}(t)}{m(t)} \right).
\end{align}

\noindent By applying the product rule and triangle inequality to $T_{\lambda(t), \gamma(t)}(x(t)) = \lambda^{-1}(t) \left( \lambda(t) T_{\lambda(t), \gamma(t)}(x(t)) \right)$, we obtain:
\begin{align}
\left\| \frac{d}{dt} T_{\lambda(t), \gamma(t)}(x(t)) \right\|_{\mathcal{H}} &\le \frac{1}{\lambda(t)} \left\| \frac{d}{dt} \left( \lambda(t) T_{\lambda(t), \gamma(t)}(x(t)) \right) \right\|_{\mathcal{H}} + \frac{|\dot{\lambda}(t)|}{\lambda^2(t)} \left\| \lambda(t) T_{\lambda(t), \gamma(t)}(x(t)) \right\|_{\mathcal{H}} \nonumber \\
&\le \frac{1}{\lambda(t)} \left( \mathcal{O}\left( \frac{\dot{m}(t)}{m(t)} \right) + \mathcal{O}\left( \frac{\dot{m}(t)}{m(t)} \right) \cdot \mathcal{O}(1) \right) \nonumber \\
&= \mathcal{O}\left( \frac{\dot{m}(t)}{m(t) \lambda(t)} \right), \label{eq:dotT_m_final_result}
\end{align}
where we have utilized \eqref{T__first_rate}, condition {\rm (C7)}, and Lemma~\ref{lem:sharp_sensitivity_final}(iii).

\noindent \textbf{(3) Improved Trajectory and Residual Rates:}
For assertion (iv), by incorporating system \eqref{eq:SHDS}, Theorem \ref{thm:pointwise_magnitudes_layered} (i)--(iii), and conditions (C8)--(C9), it follows that as $t \to +\infty$:
\begin{align}
\|\ddot{x}(t)\|_{\mathcal{H}} &\le |\alpha(t)| \|\dot{x}(t)\|_{\mathcal{H}} + |\beta(t)| \left\| \dot{T}_{\lambda(t), \gamma(t)}(x(t)) \right\|_{\mathcal{H}} + |b(t)| \|T_{\lambda(t), \gamma(t)}(x(t))\|_{\mathcal{H}} \nonumber \\
&\le \mathcal{O}\left( \frac{\dot{m}(t)}{m(t)} \right) \mathcal{O}\left( \frac{1}{m(t)} \right) + \mathcal{O}\left( \frac{m(t)}{\dot{m}(t)} \right) \mathcal{O}\left( \frac{\dot{m}(t)}{m(t)\lambda(t)} \right) + \mathcal{O}\left(1\right)\mathcal{O}\left(\frac{1}{\lambda(t)}\right)\nonumber \\
&= \mathcal{O}\left( \frac{\dot{m}(t)}{m^2(t)} \right) + \mathcal{O}\left(\frac{1}{m^2(t)}\right)  \nonumber \\
&= \mathcal{O}\left( \frac{\dot{m}(t)}{m^2(t)} \right).
\label{eq:accel_magnitude_step_O_refined}
\end{align}

Now, we focus on improving the convergence rate of $\|\dot{x}(t)\|_{\mathcal{H}}$ to validate assertion (v). Let $k_{\beta}>0$ be such that $\beta^2(t)\le k_{\beta}^2m^2(t)/\dot m^2(t)$ for all sufficiently large $t$, which follows from condition {\rm (C8)}. By the decomposition $\dot{x}(t)=v(t)-\beta(t)T_{\lambda(t),\gamma(t)}(x(t))$, Theorem~\ref{thm:m_weight_bounds}(iii)--(iv), and conditions {\rm (C10)}--{\rm (C11)}, we have
\begin{align}
\int_{t_0}^{+\infty}m(t)\dot m(t)\|\dot{x}(t)\|_{\mathcal{H}}^2dt
&\le 2\int_{t_0}^{+\infty}m(t)\dot m(t)\|v(t)\|_{\mathcal{H}}^2dt \nonumber\\
&\quad +2\int_{t_0}^{+\infty}m(t)\dot m(t)\beta^2(t)\|T_{\lambda(t),\gamma(t)}(x(t))\|_{\mathcal{H}}^2dt \nonumber\\
&\le \frac{2}{C_v}\int_{t_0}^{+\infty}\mathcal{W}_v(t)\|v(t)\|_{\mathcal{H}}^2dt \nonumber\\
&\quad +K_0\int_{t_0}^{+\infty}\mathcal{W}_T(t)\|T_{\lambda(t),\gamma(t)}(x(t))\|_{\mathcal{H}}^2dt<+\infty,
\label{eq:velocity_log_integral}
\end{align}
where $K_0:=\frac{2k_{\beta}^2}{K_1(\underline m')^2}>0$. Indeed, the last term is controlled by $\mathcal{W}_T(t)$ because condition {\rm (C11)} gives $\mathcal{W}_T(t)\ge \frac{w}{2}\frac{\dot m(t)\lambda^2(t)}{m(t)}$ for sufficient large $t$, and condition {\rm (C9)} together with $\dot m(t)\ge\underline m'>0$ yields $\mathcal{W}_T(t)\ge K_1m^3(t)\dot m(t)$ for some  $K_1:=\frac{wC_\lambda^2}{2}>0$.

Set $F(t):=m(t)\dot{x}(t)$ and $r(t):=\dot m(t)/m(t)$. Then $F$ is locally absolutely continuous. Since $\dot m(t)\ge\underline m'>0$, one has $m(t)\to+\infty$ and $\int_{t_0}^{+\infty}r(t)dt=+\infty$. The estimate \eqref{eq:velocity_log_integral} is equivalent to $\int_{t_0}^{+\infty}r(t)\|F(t)\|_{\mathcal{H}}^2dt<+\infty$.
Moreover, by the assertion {\rm (ii)} of Theorem~\ref{thm:pointwise_magnitudes_layered} and \eqref{eq:accel_magnitude_step_O_refined}, we have
\[
\|\dot F(t)\|_{\mathcal{H}}
\le \dot m(t)\|\dot x(t)\|_{\mathcal{H}}+m(t)\|\ddot x(t)\|_{\mathcal{H}}
=\mathcal{O}\left(\frac{\dot m(t)}{m(t)}\right)
=\mathcal{O}(r(t)).
\]
Thus, utilizing Lemma~\ref{lem:weighted_decay_criterion} gives $m(t)\dot{x}(t)\to0$ as $t \to +\infty$, which proves assertion (v).

For assertion (vi), given condition (C11), there exists $t_2 \ge t_0$ such that for all $t \ge t_2$, $\frac{m(t) \mathcal{W}_T(t)}{\dot{m}(t) \lambda^2(t)} \ge \frac{w}{2} > 0$. 
From Theorem \ref{thm:m_weight_bounds} (iii), it follows that
\begin{align}
\frac{w}{2} \int_{t_2}^{+\infty} \frac{\dot{m}(t)}{m(t)} \left\| \lambda(t) T_{\lambda(t), \gamma(t)}(x(t)) \right\|_{\mathcal{H}}^2 dt &\le \int_{t_2}^{+\infty} \left( \frac{m(t) \mathcal{W}_T(t)}{\dot{m}(t) \lambda^2(t)} \right) \frac{\dot{m}(t)}{m(t)} \left\| \lambda(t) T_{\lambda(t), \gamma(t)}(x(t)) \right\|_{\mathcal{H}}^2 dt \nonumber \\
&= \int_{t_2}^{+\infty} \mathcal{W}_T(t) \left\| T_{\lambda(t), \gamma(t)}(x(t)) \right\|_{\mathcal{H}}^2 dt < +\infty, 
\label{eq:U_integral_log_m}
\end{align}
which implies that $\int_{t_0}^{+\infty} \frac{\dot{m}(t)}{m(t)} \left\| \lambda(t) T_{\lambda(t), \gamma(t)}(x(t)) \right\|_{\mathcal{H}}^2 dt < +\infty$. Let $U(t):=\lambda(t)T_{\lambda(t),\gamma(t)}(x(t))$ and $r(t):=\frac{\dot m(t)}{m(t)}$, then the preceding estimate~\eqref{eq:U_integral_log_m} is precisely $\int_{t_0}^{+\infty}r(t)\|U(t)\|_{\mathcal{H}}^2dt<+\infty$. Furthermore, as shown in the proof of assertion {\rm (iii)}, $U$ is locally absolutely continuous and $\|\dot U(t)\|_{\mathcal{H}}=\mathcal{O}(r(t))$. Therefore, applying Lemma~\ref{lem:weighted_decay_criterion} to $U(t)$ yields $U(t)\to0$, namely $\lambda(t)\|T_{\lambda(t),\gamma(t)}(x(t))\|_{\mathcal{H}}\to0$ as $t \to +\infty$. This validates assertion (vi).

To prove the remain results, invoking the rate $\|T_{\lambda(t), \gamma(t)}(x(t))\|_{\mathcal{H}} = o(\frac{1}{\lambda(t)})$ established in Theorem \ref{thm:pointwise_magnitudes_layered} (vi), we get
\begin{align*}
\left\| A_{\gamma(t)} \left( x(t) - \gamma(t) Bx(t) \right) + Bx(t) \right\|_{\mathcal{H}} &= \frac{\lambda(t)}{\gamma(t)} \left\| T_{\lambda(t), \gamma(t)}(x(t)) \right\|_{\mathcal{H}}  = o\left( \frac{1}{\gamma(t)} \right),
\label{eq:operator_sum_magnitude_result}
\end{align*}
which establishes assertion (vii). 

\noindent By applying triangle inequality to the time derivative identity of the regularized operator sum and utilizing the results {\rm (iii)} and {\rm (vi)} from Theorem \ref{thm:pointwise_magnitudes_layered}, we verify assertion {\rm (viii)} through the following estimation:
\begin{equation*}
\begin{aligned}
\left\|\frac{d}{dt}\left(A_{\gamma(t)}\left(x(t)-\gamma(t)Bx(t)\right)+Bx(t)\right)\right\|_{\mathcal{H}} 
&\le \frac{\lambda(t)}{\gamma^2(t)} \left| \frac{d}{dt} \left( \frac{\gamma(t)}{\lambda(t)} \right) \right| \left\| \lambda(t) T_{\lambda(t), \gamma(t)}(x(t)) \right\|_{\mathcal{H}} + \frac{\lambda(t)}{\gamma(t)} \left\| \dot{T}_{\lambda(t), \gamma(t)}(x(t)) \right\|_{\mathcal{H}} \\
&\le o\left( \frac{\lambda(t) \left| \frac{d}{dt} \left( \frac{\gamma(t)}{\lambda(t)} \right) \right|}{\gamma^2(t)} \right) + \frac{\lambda(t)}{\gamma(t)} \mathcal{O}\left( \frac{\dot{m}(t)}{m(t) \lambda(t)} \right) \\
&= \mathcal{O}\left( \frac{\dot{m}(t)}{m(t) \gamma(t)} \right) + o\left( \frac{\lambda(t) \left| \frac{d}{dt} \left( \frac{\gamma(t)}{\lambda(t)} \right) \right|}{\gamma^2(t)} \right).
\end{aligned}
\end{equation*}
\end{proof}

\end{document}
